% Expository adaptation, October 7, 2026. Original mathematics: OpenAI, September 23, 2026.
% Prose revised and explanations expanded; see README.md for the source and scope.
\section{Representing a row by a default and exceptions}
\label{sec:buf-compression}

The averaging schedule describes a value at every retained destination, but we cannot keep a growing list of all destinations encountered. After each group we therefore replace the row by a common default value and a controlled number of exceptions. The list remains implicit: its entries will be located by recomputation.

We do not need a uniformly accurate raw value at every position. The algorithm subsequently applies bounded row functions, scores, and support witnesses to the row. These functions ignore sufficiently small entries, become constant beyond certain ranges, and shut off when a row is too crowded. The compression analysis uses precisely these properties.

The proof distinguishes the first group at which the allowed number of exceptions forces the compression threshold above its prescribed floor. Before that event, compression changes only very small values. At that event, either the true row has unusually many sampled large entries, or earlier errors have made small entries appear large. The first case is rare even after requiring a queried destination to be retained; the second is charged to those earlier errors. A comparison with the true row then controls subsequent updates.

This distinction is also why the mean error and averaging noise were kept separate. The mean discrepancy at a fixed destination is deterministic and can be multiplied by the probability of a rare sampling event. The noise can depend on that event and needs its own unconditional bound.

We first fix the increments supplied by exact endpoint grouping and prove
the comparison without range clamps.  We then show that clamping preserves
the relevant row functions and assemble the exact-grouping error reserve.
Section~\ref{sec:fp-blocks} compares the actual and clean calculations with
those clamps in place.

\subsection{The row properties that compression uses}
\label{subsec:buf-interface}
Fix a row task, its start rate $u$, and its first-channel endpoint rate $a$.
Write $d_y\geq0$ for its true first raw input, in units of $a$.  Its second
raw input, when present, is $v_y=\zeta\mathbf1_{\{Y_t(y)=1\}}q_y$.
Here $q_y$ is the deterministic conditional mean, $0\leq q_y\leq1$, and
$q_y=0$ whenever $d_y$ is below the positive-start threshold of the
preparatory function $h$.  Thus $d_y$ and $q_y$ are deterministic at each
specified endpoint; $v_y$ carries the prescribed narrower membership mask.
There is a fixed constant $C_{\mathrm{mass}}$ such that
\begin{equation}
 a\sum_y d_y\leq C_{\mathrm{mass}}.
 \label{eq:buf-mass}
\end{equation}
The endpoint sets in this section are the retained sets, with the stipulated
common padding on their hash arguments.

We apply a fixed finite family of pre-incoming row functions to a raw array.
Denote this family by $\mathcal H$.  It includes the main pre-functions,
auxiliary scores, and witnesses required by the preceding construction.
Each function uses the first entry, possibly the second entry, and the
$K$th largest first entry.  For each member $H\in\mathcal H$ there are fixed
positive thresholds $g_H,e_H$, with $3g_H<e_H$, such that its row gate is
zero when that rank is at least $g_H$, and its entry factor is zero when the
first entry is at most $e_H$.  The gate is constant and full in a neighborhood
of zero.  The $s$th roots of the output functions are bounded and Lipschitz
in their relevant entry and rank arguments, where $s=8$.  All these arguments
can be clipped to fixed sensitivity ranges without changing the functions.
These are exactly the threshold buffers and saturations imposed when
constructing the atoms; the thresholds for the required functions are
listed in Table~\ref{tab:atom-compression-thresholds}.  In particular,
for every real raw array, including
arrays with negative entries,
\begin{equation}
 \sum_{H\in\mathcal H}\#\{y:H(y)\ne0\}\leq C_{\mathrm{supp}},
 \label{eq:buf-row-support}
\end{equation}
where $C_{\mathrm{supp}}$ is fixed.  Indeed a nonzero gate implies that fewer
than $K$ entries exceed its close threshold, and every positive output entry
exceeds the larger entry threshold.  This observation applies also to the
broader witnesses.

Choose $\eta>0$ below all positive entry thresholds, all positive-start
thresholds for rank sensitivity, and the threshold for $h$, with enough
fixed separation that perturbations of size several times $\eta$ remain
inside these flat regions.  Choose $\zeta\leq\eta/100$.  These constants are
fixed throughout the proof.

Let $d_1>d_2>\cdots>d_{m_*}=1$ be the decreasing powers of two indexing
slack groups.  Empty groups are included.  Let $S_i=(S_i^1,S_i^2)$ be the
uncompressed prefix through group $i$, and let
$\Delta_i=S_i-S_{i-1}$, with $S_0=0$.  Define the clipped endpoint errors
and their row supremum by
\begin{equation}
 \begin{aligned}
 e_i(y)&=\max\left\{
 \min\left(1,\frac{|S_i^1(y)-d_y|}{1+d_y}\right),
 \min(1,|S_i^2(y)-v_y|)\right\},\\
 e_i&=\sup_y e_i(y).
 \end{aligned}
 \label{eq:buf-errors}
\end{equation}
omitting channel two when absent.  Write $\beta_i(y)$ for the maximum of
the mean discrepancies, with the same clipping and first-channel
normalization as in Equation~\eqref{eq:buf-errors}, and $\gamma_i(y)$ for
the maximum absolute prefix noise.  Thus
$e_i(y)\leq\beta_i(y)+\gamma_i(y)$.  At a specified row and endpoint, the
first-channel mean discrepancy is deterministic.  The second-channel mean
discrepancy is a deterministic magnitude times its narrower membership
mask.  Prefix noise may depend on the shared environment.
The deterministic assertion is a hypothesis of the compression argument:
the slice average defining a mean must vary every environment coordinate
on which its sample function depends.  Lemma~\ref{lem:est-prefix} states
this requirement, and Lemma~\ref{lem:fp-clean} verifies it for the clean
calculations to which we will apply the result.  We will use conditional
rarity to control these fixed mean-error magnitudes; the noise will be
controlled by its unconditional norm bound.

For row integration use
\[
 \int F\,d\mu=\frac1N\mathbb E\sum_x\frac{X_u(x)}u F(x).
\]
Endpoint integration means inserting the sum over retained endpoints in
this expression.  The prefix estimates used as hypotheses below are
\begin{equation}
 \left(\int\sum_y\beta_i(y)^s\,d\mu\right)^{1/s}
       \leq c_*2^{-p+B_1d_i},\qquad
 \left(\int\sum_y\gamma_i(y)^s\,d\mu\right)^{1/s}
       \leq c_*2^{-p-K_nd_i}.
 \label{eq:buf-prefix-interface}
\end{equation}
Lemma~\ref{lem:est-prefix} supplies these estimates for the additive
schedule.  Here and below endpoint sums use the relevant masks.  In
particular these bounds control the $L_s(\mu)$ norms of the corresponding
row suprema, which we denote by $\beta_i$ and $\gamma_i$.

\subsection{The compression update}

The update needs to preserve large entries, suppress small ones, and remain stable when its inputs change slightly. Hard deletion at a threshold would fail the last requirement. We instead clip the displacement from a new default, with a clipping radius that grows continuously from zero.

Three quantities determine the threshold: the previous default, a prescribed positive floor, and an order statistic that enforces the list capacity. The additional $K$th-rank adjustment makes the default compatible with the row gates. The formulas below specify this adjustment; Lemma~\ref{lem:buf-update} isolates its elementary properties before the error analysis begins.

Adjoin arbitrarily many dummy positions, all with target and increments
zero.  A compressed array has default $(z,0)$, shared also by the dummies,
and a finite list of exceptions.  Initially $z=0$ and the list is empty.
Order statistics of an array with this default mean the order statistics
of its finite exceptional values with as many default values as needed.
For a rank of order $r$, the finite list together with at least $r$
default copies gives exactly the prescribed rank.  Additional real list
positions whose values equal the default do not change this rank.
Dummy positions serve only this rank computation; endpoint sums in the
error bounds are over real retained IDs.  The margin statement of
Lemma~\ref{lem:atom-margins} ensures that a default position cannot expose
a row port.
At group $j$ add $\Delta_j$ by actual endpoint equality, obtaining $w_y$,
and set
\begin{equation}
 \begin{aligned}
 \rho_j&=2^{-c_1d_j},\qquad n_j=2^{c_2d_j},\qquad
 m_y=\max_\nu|w_y^\nu|,\\
 u_0&=\max\{z,\rho_j,\operatorname{rank}_{n_j}(m)\},\\
 q_0&=\max\{0,\operatorname{rank}_K(w^1)\},\\
 z'&=\max\{u_0,\min(q_0,3u_0)\}.
 \end{aligned}
 \label{eq:buf-thresholds}
\end{equation}
Here $c_1,c_2$ are fixed sufficiently large positive integers, and
$c_2$ is chosen also so that $n_j\geq K$ in every group.
With $b=(z',0)$, the updated array is
\begin{equation}
 T_j^\nu(y)=b^\nu+
 \operatorname{clip}\bigl(w_y^\nu-b^\nu,
                \pm7(m_y-2u_0)_+\bigr).
 \label{eq:buf-update}
\end{equation}
The notation $\operatorname{clip}(r,\pm A)=\max(-A,\min(r,A))$ is used.

\begin{lemma}[Local compression properties]
\label{lem:buf-update}
The update has the following properties.
It replaces every position of size at most $2u_0$ by the new default, is
identity at every position of size at least $3u_0$, and lies coordinatewise
between the input and the new default.  The default is nondecreasing.
The map from the input array and old default to the updated array and new
default is uniformly Lipschitz in the sup norm.  A list containing all
exceptions can be chosen with at most $n_j$ positions, using only the
previous list and the new candidates.
\end{lemma}
\begin{proof}
First inspect the two ranges where the update is especially simple. At size at most $2u_0$, its clipping radius is zero, so the entry becomes the default.  Since $z'\leq3u_0$, if $m_y\geq3u_0$
then
\[
 |w_y^1-z'|\leq m_y+3u_0\leq7(m_y-2u_0),
 \qquad |w_y^2|\leq7(m_y-2u_0),
\]
which proves identity there.  Clipping a displacement toward zero proves
the coordinatewise assertion.  Also $z'\geq u_0\geq z$.

If two input arrays and their defaults differ by at most $\delta$, their
size arrays, size ranks, $u_0$'s, and $q_0$'s differ by at most $\delta$.
Their new defaults differ by at most $3\delta$, and their clipping radii
by at most $21\delta$.  The elementary bound
\[
 |\operatorname{clip}(r,\pm A)-\operatorname{clip}(r',\pm A')|
       \leq |r-r'|+|A-A'|
\]
therefore gives, for example, a Lipschitz constant 28 for the full update.
This constant is independent of the number and magnitudes of entries.

Fewer than $n_j$ positions have size strictly larger than $u_0$.
A coarse test supported above $u_0$ and full above $2u_0$ retains all
possible exceptions within the stated cap.  A position outside the old
list and new candidate set has input $(z,0)$, has size at most $u_0$,
and hence becomes the new default.  No other position need be inspected.
The selected list may retain some real positions that become exactly
default after compression.  Such positions still represent the same
full array.  No padding dummy is selected, since its input size is at
most $u_0$.
\end{proof}

In particular an absolute perturbation introduced at group $i$ propagates
through the remaining updates and the final root functions with factor
at most
\begin{equation}
 C C_1^{m_*-i+1}
       \leq C(d_i+1)^{\log_2 C_1},
 \label{eq:buf-propagation}
\end{equation}
for a fixed $C_1>1$, enlarged below as needed.  This statement concerns
full arrays and defaults; it does not require their exceptional lists to
agree.

\subsection{Compare with a row whose relevant outputs stay fixed}

Before integrating over randomness, we establish deterministic stability statements for one row. The most useful reference is the true row together with the small changes already caused by early compression. From the first overflow onward, run compression on this reference without adding further increments. We will show that its relevant row functions stay fixed. The actual calculation can then be compared with this reference by the prefix-error bounds.

Call group $j$ an overflow if its $u_0$ exceeds $\rho_j$.  Before the first
overflow every compression change has magnitude at most $6\rho_i$ and
occurs only at input size below $3\rho_i$.  Because successive slacks
halve, choosing $c_1$ large gives, uniformly in $j$,
\begin{equation}
 6\sum_{i<j}\rho_i\leq\rho_j/100,
 \qquad 6\sum_i\rho_i\leq\eta/100.
 \label{eq:buf-small-floors}
\end{equation}
The same estimates cover dummy defaults.  For example, before the first
overflow the preceding default is at most $3\rho_{j-1}\ll\rho_j$.

\begin{lemma}[Preloss and stationary reference]
\label{lem:buf-reference}
Let $L_{j-1}=T_{j-1}-S_{j-1}$ before the first overflow at $j$.
Then $\|L_{j-1}\|_\infty\leq6\sum_{i<j}\rho_i$, and at every endpoint
with $d_y\geq\eta$,
\begin{equation}
 \max_\nu|L_{j-1}^\nu(y)|
       \leq C\sum_{i<j}\rho_i e_i(y).
 \label{eq:buf-preloss}
\end{equation}
Start a reference calculation at the input of group $j$ with the true
array $(d,v)+L_{j-1}$, the same preceding default, and no further
increments.  Every subsequent compression preserves its relevant row
outputs.  Their root discrepancy from the true outputs is bounded by
$C\sum_{i<j}\rho_i e_i$.
\end{lemma}
\begin{proof}
First account for the changes before the first overflow. Their uniform bound is small, but at a truly substantial entry we need a stronger bound that can be charged to earlier estimation error. The uniform preloss bound follows by adding the changes described above.
If a change occurs at group $i$ at a position with $d_y\geq\eta$, then
$|w_y^1|\leq3\rho_i$ and the preceding preloss is much smaller than
$\rho_i$.  Thus
\[
 \frac{|S_i^1(y)-d_y|}{1+d_y}
 \geq \frac{d_y-3\rho_i-6\sum_{k<i}\rho_k}{1+d_y}
 \geq c_\eta>0.
\]
The change of size at most $6\rho_i$ is consequently bounded by
$C\rho_i e_i(y)$.  Adding such changes proves
Equation~\eqref{eq:buf-preloss}.

Next show that repeated compression without new increments leaves the reference outputs unchanged. We track its $K$th-largest first entry, because that is the statistic used by the row gates. In the reference array negative first magnitudes and all second magnitudes
are at most $\eta$, by the choices of $\zeta$ and the floors.
Let $q$ be its initial $K$th first-coordinate rank.  The default is
nonnegative, so $q\geq0$.  If $q>\eta$, then $n_i\geq K$ implies that the
$n_i$th size rank is at most $q$.  The floor and preceding default are
also at most $q$, so $u_0\leq q$ and $z'\leq q$.  Entries at most $q$
stay at most $q$.  Entries at least $q$ are unchanged if $q\geq3u_0$;
otherwise $z'=q$, and they stay at least $q$.  The $K$th rank is therefore
unchanged.  Negative first magnitudes and second magnitudes do not grow,
so this argument repeats at every subsequent group.  If instead $q\leq\eta$,
the size rank, $u_0$, and new default are at most $\eta$, and the next
$K$th rank is again at most $\eta$.

Consider an individual row function.  In the first case its rank stays
fixed.  If its gate is not already zero, $q<g_H$, so every entry in its
support has first coordinate greater than $e_H>3q\geq3u_0$ and both
channels there are untouched.  A position below entry support cannot
cross into support because the default is at most $q<e_H$.
In the second case the rank stays in the common full-gate region; the
same identity argument applies to every position in entry support.
Thus each reference output is preserved.

Finally, positions with $d_y<\eta$ stay below every entry threshold and
cannot change a rank within its sensitivity region.  Perturbations at
all other positions satisfy Equation~\eqref{eq:buf-preloss}.  Clip first
entries and ranks to the fixed sensitivity ranges and use root-Lipschitz
continuity of the row functions.  This proves the last assertion.
\end{proof}

If no overflow ever occurs, the same proof, together with the final prefix
error, gives the bound
\begin{equation}
 \sup_{H,y}|H(T_{m_*};y)^{1/s}-H(d,v;y)^{1/s}|
       \leq C\left(e_{m_*}+\sum_i\rho_i e_i\right).
 \label{eq:buf-no-overflow}
\end{equation}
For completeness, if the expression on the right is not small, bounded
output ranges give the assertion after enlarging $C$.  Otherwise low-target
positions remain below sensitivity.  At other positions clipped first
entry discrepancies are bounded by a constant times their relative error,
and second entry discrepancies by their absolute error.  Adding the
preloss estimate and applying the root-Lipschitz bounds proves the claim.

A raw first-channel error is measured relative to $1+d_y$, so it can be large in absolute value when the true entry is large. The next lemma must therefore use saturation and the row gates, rather than merely multiply an absolute sup-norm error by a Lipschitz constant.

\begin{lemma}[Stability after the first overflow]
\label{lem:buf-relative-stability}
At a first overflow in group $j$, the discrepancy of final row-root outputs
from the stationary reference of Lemma~\ref{lem:buf-reference} is at most
\begin{equation}
 C\sum_{i\geq j} C_1^{m_*-i+1}e_i.
 \label{eq:buf-relative-propagation}
\end{equation}
The constant does not depend on the largest true first entry.
\end{lemma}
\begin{proof}
Write $E$ for the sum in Equation~\eqref{eq:buf-relative-propagation}
without the initial constant.  If $E$ exceeds a sufficiently small fixed
constant, bounded output ranges suffice.  Otherwise all clipped errors
involved are strictly below 1.  Hence
\[
 |S_i^1(y)-d_y|\leq(1+d_y)e_i,
 \qquad |S_i^2(y)-v_y|\leq e_i,
\]
and later increments have first magnitude at most
$(1+d_y)(e_i+e_{i-1})$ and second magnitude at most $e_i+e_{i-1}$.
The input difference at group $j$ is precisely $S_j-(d,v)$; both
calculations include the same previous preloss.
The factors $1+d_y$ prevent a direct absolute-error bound with a uniform
constant when a true entry is large.  The following two cases use the
row gate and saturation to remove that dependence.

Choose $G$ larger than every gate close threshold.  A baseline at least
$G$ closes all gates, since the dummy entries force every first-coordinate
rank to be at least the baseline.  Closure is permanent under the
unclamped update.  Choose $L_0$ much larger than $G$, all sensitivity
thresholds, and 1.  If at least $K$ true first entries exceed $L_0$, the
reference gates are closed.  In the perturbed calculation either some
$u_0\geq G$, which forces permanent closure, or every $u_0<G$.
In the latter case those $K$ coordinates stay larger than half their
true value: their initial error and the total later additive variation
are at most $C(1+d_y)\sum_{i\geq j}e_i$, and their first value is then
larger than $3G$.  Inductively each such position lies in the compression
identity region.  These $K$ positive entries close all gates as well.

It remains to consider the case of fewer than $K$ true entries above
$L_0$.  Lemma~\ref{lem:buf-reference} bounds every reference rank and
threshold by $2L_0$.  Fix $L_1$ much larger than $L_0$ and all upper
saturation thresholds.  At the at most $K-1$ positions where $d_y>L_1$,
replace the first input at group $j$, in both calculations, by one common
positive sentinel larger than $L_1$, and give that coordinate zero later
first-channel increments.  Do not modify second channels.
The reference ranks, thresholds, and outputs are unchanged: these
coordinates stay above the relevant ranks and in all upper saturation
regions.
The sentinels are comparison devices in this proof.  They need not be
identified or stored by the compression algorithm.  Fewer than $K$ such
coordinates occur, so replacing them by common large values removes
their possibly large absolute differences while preserving every rank
that lies in a sensitivity range.

All remaining first-coordinate input errors and increments now have
absolute bounds by $C(1+L_1)e_i$ and
$C(1+L_1)(e_i+e_{i-1})$, respectively.  Lemma~\ref{lem:buf-update},
iterated from $j$ to the last group, therefore bounds every array and
threshold discrepancy by $CE$, on taking $C_1$ larger than the fixed
update constant.  If $E$ is sufficiently small, all these thresholds
remain below $2L_0+1$.  The final root-output discrepancy is at most $CE$.

To remove the modification in the perturbed calculation, argue inductively.
Every original coordinate with $d_y>L_1$ remains above $d_y/2$, because
its total raw additive variation is at most $C(1+d_y)\sum e_i$.
It is thus above the bounded thresholds and in the identity and saturation
regions.  Replacing fewer than $K$ such positive coordinates by sentinels
cannot change either a $K$th first rank or an $n_i$th size rank whose
sentinel-computed value is at most $2L_0+1$: both the original and modified
values are above that bound, and $n_i\geq K$.
The defaults and thresholds therefore agree, the compression is identity
in both channels at these positions, and all other coordinates agree.
The relevant output behavior is the same.  This completes the induction
and the proof.
\end{proof}

\subsection{Distinguish truly crowded rows from estimation errors}

The deterministic comparisons leave two cases to integrate: no overflow,
and the first overflow.  To handle the latter, we separate a large retained
set of true heavy endpoints from overflow caused by errors at light
endpoints.
For this subsection let $\mathcal O_j$ denote the event that group $j$
is the first overflow.  These events are pairwise disjoint.  The
heavy-count event defined next depends only on the true row and its
retained endpoints; it does not itself assert that any overflow occurs.

At first overflow in group $j$, the previous default is smaller than
$\rho_j$.  Thus at least $n_j$ retained real positions have input size
strictly greater than $\rho_j$.  Define the true heavy-count event
\[
 \Theta_j=\left\{\#\{y\text{ retained}:d_y>\rho_j/4\}\geq n_j/2\right\}.
\]
The dummy positions do not contribute to this count, since their input
is the preceding default.  On $\mathcal O_j\cap\Theta_j^c$, at least
$n_j/2$ of the overflowing positions have
$d_y\leq\rho_j/4$.  Their second targets vanish, because the floors lie
below the $h$-threshold.  Equation~\eqref{eq:buf-small-floors} then implies
$e_j(y)\geq c\rho_j$ at all these positions.  Consequently the indicator
of this false-type first overflow is bounded by
\begin{equation}
 \frac{C}{n_j\rho_j^s}\sum_y e_j(y)^s.
 \label{eq:buf-false-charge}
\end{equation}
A bounded row-output error may be charged to this indicator.
On $\mathcal O_j\cap\Theta_j$, we instead use the stability estimate
and the conditional rarity of $\Theta_j$ at each endpoint.

\begin{lemma}[Conditional rarity of a true heavy count]
\label{lem:buf-rarity}
At every specified endpoint,
\begin{equation}
 \mathbb P(\Theta_j\mid X_u(x)=Y_a(y)=1)
      \leq\frac{C}{n_j\rho_j}.
 \label{eq:buf-rarity}
\end{equation}
The same bound holds with the narrower endpoint condition $Y_t(y)=1$.
Moreover, for every prefix mean row supremum,
\begin{equation}
 \|\mathbf1_{\Theta_j}\beta_i\|_{L_s(\mu)}
      \leq Cc_*2^{-p+B_1d_i}(n_j\rho_j)^{-1/s}.
 \label{eq:buf-masked-mean}
\end{equation}
\end{lemma}
\begin{proof}
Bound the number of genuinely large entries using row mass, then bound how many can be retained after one destination is forced to survive. By Equation~\eqref{eq:buf-mass}, the deterministic heavy set
\[
 A_j(x)=\{y:d_y>\rho_j/4\}
 \quad\text{satisfies}\quad
 |A_j(x)|\leq\frac{4C_{\mathrm{mass}}}{a\rho_j}.
\]
Conditional on one specified endpoint
hit, each other endpoint in this set is retained with probability at most
$Ca$, by Lemma~\ref{lem:env-ranks}.  This remains true if the
specified hit uses the narrower rate.  The forced endpoint contributes
at most 1, so the expected retained heavy count is at most
$1+Ca|A_j(x)|\leq C'/\rho_j$.
The independent start condition has no effect.  Markov's inequality
at the threshold $n_j/2$ proves Equation~\eqref{eq:buf-rarity}.
This calculation uses the conditional pair bound and linearity of
expectation; it permits dependence among the heavy-endpoint hits.

To exploit rarity in the mean-error estimate, condition separately at each destination before summing. This works because the magnitude of its mean error is fixed under that conditioning. For the first-channel mean error at a fixed row and endpoint, write its
deterministic magnitude as $b_i(y)$.  Then
\[
 \mathbf1_{\Theta_j}\sup_{y:Y_a(y)=1}b_i(y)^s
       \leq\sum_y\mathbf1_{\Theta_j}Y_a(y)b_i(y)^s.
\]
Condition on the start and that endpoint before taking expectations.
Equation~\eqref{eq:buf-rarity} bounds the resulting expectation by
$C/(n_j\rho_j)$ times the endpoint-integrated mean-error power.
Indeed the factor $b_i(y)^s$ is fixed under this conditioning, so
\[
 \mathbb E\bigl[\mathbf1_{\Theta_j}X_u(x)Y_a(y)b_i(y)^s\bigr]
 \leq\frac{C}{n_j\rho_j}
       \mathbb E\bigl[X_u(x)Y_a(y)b_i(y)^s\bigr].
\]
The second-channel mean error has a deterministic magnitude multiplied
by $Y_t(y)$; repeat the same argument using the narrower conditional
bound.  Taking roots and using Equation~\eqref{eq:buf-prefix-interface}
proves Equation~\eqref{eq:buf-masked-mean}.  The conditioning step uses
the deterministic mean-error magnitude at the specified endpoint.
Prefix noise can depend on the heavy-count event, so no analogous
conditional estimate is asserted for it.
\end{proof}

\begin{proposition}[Integrated compression bound]
\label{prop:buf-error}
For fixed sufficiently large $c_1$, then $c_2$, and sufficiently rapid
prefix-noise decay $K_n$, compression preserves the prescribed row error
budget, with an arbitrarily small fixed coefficient when $c_*$ is chosen
small.  More precisely,
\begin{equation}
 \left(\int\sum_{H\in\mathcal H}\sum_y
  |H(T_{m_*};y)^{1/s}-H(d,v;y)^{1/s}|^s\,d\mu\right)^{1/s}
       \leq C_{\mathrm{buf}}c_*2^{-p},
 \label{eq:buf-integrated}
\end{equation}
where $C_{\mathrm{buf}}$ is fixed, independent of the number of groups,
row width, and requested precision.
\end{proposition}
\begin{proof}
We add four contributions: the small changes preceding overflow, the final prefix error when no overflow occurs, overflow caused by erroneous large entries, and overflow in a truly crowded row. In the last case, mean error and noise are treated separately.

Equation~\eqref{eq:buf-row-support} converts every row sup bound above
into a bound for its summed output power with a fixed factor.
In fact the support of a discrepancy lies in the union of the true and
computed supports, whose total size over $\mathcal H$ is at most
$2C_{\mathrm{supp}}$.  Thus this step introduces no factor from the
ambient number of endpoints.
The preloss contributions have norm at most
\[
 Cc_*2^{-p}\sum_i\rho_i
             \bigl(2^{B_1d_i}+2^{-K_nd_i}\bigr).
\]
On the event of no overflow, Equation~\eqref{eq:buf-no-overflow} adds the
final-prefix cost, at most $Cc_*2^{-p}(2^{B_1}+2^{-K_n})$.
By Equation~\eqref{eq:buf-false-charge} and the disjointness of the
possible first-overflow events, the total false-type error power is at
most
\begin{equation}
 Cc_*^s2^{-sp}\sum_j
       (n_j\rho_j^s)^{-1}
       \bigl(2^{sB_1d_j}+2^{-sK_nd_j}\bigr).
 \label{eq:buf-false-sum}
\end{equation}

For a true-type first overflow at $j$, use
Lemma~\ref{lem:buf-relative-stability} and split $e_i\leq\beta_i+\gamma_i$.
For the mean part, enlarge the actual first-overflow event to $\Theta_j$,
apply Equation~\eqref{eq:buf-masked-mean}, and use Minkowski's inequality.
Its total norm is bounded by
\begin{equation}
 Cc_*2^{-p}\sum_j(n_j\rho_j)^{-1/s}
       \sum_{i\geq j}C_1^{m_*-i+1}2^{B_1d_i}.
 \label{eq:buf-true-sum}
\end{equation}
For the noise part, use the uniqueness of the first-overflow group on
each row.  Its tail noise sum is bounded pointwise by the single full
sum $\sum_iC_1^{m_*-i+1}\gamma_i$, regardless of dependence between
noise and overflow.  More explicitly,
\[
 \sum_j\mathbf1_{\mathcal O_j}
       \sum_{i\geq j}C_1^{m_*-i+1}\gamma_i
 \leq\sum_i C_1^{m_*-i+1}\gamma_i,
\]
because at most one indicator on the left is nonzero.  The absolute
prefix-noise bound therefore gives
norm at most
\begin{equation}
 Cc_*2^{-p}\sum_iC_1^{m_*-i+1}2^{-K_nd_i},
 \label{eq:buf-noise-sum}
\end{equation}
without an extra sum over $j$.

Choose $c_1>B_1$ large enough for the floor conditions.  Choose
$c_2>sc_1+sB_1$ with a fixed positive margin and also large enough that
$n_j\geq K$.  Equation~\eqref{eq:buf-false-sum} then converges uniformly.
For its mean term the exponential factor is exactly
$2^{-(c_2-sc_1-sB_1)d_j}$.
For Equation~\eqref{eq:buf-true-sum}, its inner sum is at most
$C(d_j+1)^A2^{B_1d_j}$ for a fixed $A$, by
Equation~\eqref{eq:buf-propagation} and dyadic grouping; the stronger
choice of $c_2$ just made ensures convergence of the outer sum, whose
exponential factor is
$2^{-((c_2-c_1)/s-B_1)d_j}$.
Taking $K_n$ sufficiently large handles Equation~\eqref{eq:buf-noise-sum}.
The preloss sum converges by $c_1>B_1$.  All constants are fixed before
choosing the arbitrarily small prefix coefficient $c_*$, giving
Equation~\eqref{eq:buf-integrated} and the asserted spare coefficient.
\end{proof}

Proposition~\ref{prop:buf-error} bounds the compressed pre-incoming
functions, including scores and witnesses.  Before assembling the final
atom estimates, we impose the numerical range bounds needed to evaluate
these arrays.

\subsection{Bound numerical ranges without changing the needed outputs}

The algorithm must also avoid arbitrarily large numerical values. A clamp can be harmless if its threshold lies beyond every sensitivity range and beyond all the remaining additive variation: a value cut off there cannot later return to a range where its exact magnitude matters. This is the tail-margin argument below.

In exact endpoint grouping, the additive increments are independent of
which exceptions are stored.  The path description and bounded child
ranges give a fixed $D$ for which
$\|\Delta_i\|_\infty\leq2^{Dd_i}$, after increasing $D$ to include
immediate term counts.  This also bounds a group when all its same-key
increments are summed into one coordinate, with a possibly larger fixed
$D$.  Fix a common gate closure level $G\ge\max_H g_H$, and then
choose $L>10G$ above every sensitivity threshold.  Let
$R_i=2^{B_2d_i}$ and choose $B_2$ sufficiently large that
\begin{equation}
 R_i\geq 10L+10\sum_{k>i}\|\Delta_k\|_\infty
 \quad\text{for every group }i.
 \label{eq:buf-clamp-margin}
\end{equation}
Dyadic grouping permits this uniform choice because $d_{i+1}=d_i/2$.
After each compression, clip every coordinate and the baseline
componentwise to $[-R_i,R_i]$.
An off-list coordinate is clipped exactly as the baseline is, so it
remains a default coordinate.  These clamps therefore preserve the
finite-list representation.

\begin{lemma}[Range clamping]
\label{lem:buf-clamping}
For fixed common increments satisfying
Equation~\eqref{eq:buf-clamp-margin}, these clamps do not change the
final relevant row functions.  Each clamp is sup-norm nonexpansive, so
in comparisons of two calculations it preserves the uniform Lipschitz
propagation bound.  The analogous assertion holds for vector sums clipped
beyond their final sensitivity range and all remaining additive variation.
\end{lemma}
\begin{proof}
Compare the clamped and unclamped calculations with the same increments.
For the fixed level $G<L/10$, once both baselines are at
least $G$, all gates are zero because of the dummies, and remain so:
compression cannot lower the baseline and every clamp bound exceeds $G$.
It remains to compare the calculations before this happens.

Maintain the following invariant.  Their baselines agree while below $G$.
Their coordinates either agree, or have matching signs and magnitudes
well beyond $L$; any such differing component has enough remaining
margin that the total future increments cannot bring it into
$[-L,L]$.  At the first clamp of a component this follows from
Equation~\eqref{eq:buf-clamp-margin}.  A subsequent clamp restores the
same margin using the current $R_i$.  Between clamps, common additions
consume at most the available tail margin.

After an addition, entries whose components differ therefore have size
larger than $L$ in both calculations.  Their size-rank statistics either
are equal below $G$ or are both at least $G$: differing values lie above
this comparison range and all other size values agree.
In the latter case both new baselines close the gates.  In the former
case the two values of $u_0$ agree and are below $G$.
Only the first-coordinate rank clipped at $3u_0$ is needed to compute
$z'$.  Differing positive first components are above this range;
differing negative first components are below zero and do not affect
its nonnegative clipped value.  Consequently the two new baselines
agree as well.  If this baseline is at least $G$, both calculations
close.  Otherwise every differing position lies in the compression
identity region because its size exceeds $L>3G$; equal coordinates at
other positions undergo the same update.  Thus the invariant continues.

Before closure, all rank arguments and all entry arguments within their
sensitivity ranges agree.  Differing large components have matching
signs and identical saturated behavior.  The final row functions
therefore agree.  A scalar clipping map is 1-Lipschitz, which gives the
comparison assertion.  For a vector sum, the same tail-margin argument
ensures that a clipped value stays on the same saturated side until the
final bounded clipping; no rank argument is needed.
\end{proof}

The equality assertion of Lemma~\ref{lem:buf-clamping} is used for the
clean error bound: exact endpoint grouping supplies the same increment
arrays regardless of the chosen list presentation.  It preserves the
values of every relevant $H\in\mathcal H$, including scores and smooth
witnesses.  It makes no assertion that finite-accuracy access tests or
selected descriptors agree with those of the unclamped calculation.
In the adaptive comparison of Section~\ref{sec:fp-blocks}, both the actual
and clean calculations are clamped from their definition.  Their
possibly different bucket increments are compared directly, and the
nonexpansiveness part of Lemma~\ref{lem:buf-clamping} preserves the local
perturbation bound.  This use allows different real representatives,
including retained positions whose current value equals the default.

\subsection{The accuracy available before fingerprints}

The prefix estimates and compression bound now give the current task's
error estimate under exact endpoint grouping.  The hypotheses still
require each conditional average to vary every random coordinate used
by its sample function.  We combine the compressed pre-incoming bounds
with the atom assembly comparison; the next section will verify this
averaging hypothesis for the block-clean calculation.

We now combine the averaging bounds, compression stability, and table assembly. The coefficient is chosen to use only half the available error bound. The remaining half will pay for replacing exact destination equality by short fingerprints. This is an allocation of constants within the same estimator, not an additional approximation assumption.

\begin{proposition}[Error reserve with exact endpoint grouping]\label{prop:est-clean}
Suppose the earlier row and vector tasks have error bound $C_b2^{-p'}$
and the support and common-parameter properties of the atom section.
Use exact comparisons locally, and vary every variable on which the
sample terms depend.  After compression and the range clamps above, all
active row outputs, pre-incoming scores, and witnesses have error at most
$(C_b/2)2^{-p}$.  The same holds for the final clipped vectors and the
assembled correction outputs, after choosing the fixed constants in the
order specified below.
\end{proposition}
\begin{proof}
To apply Proposition~\ref{prop:buf-error}, use the deterministic-mean and
absolute-noise bounds of Lemma~\ref{lem:est-prefix}.  The first target
is nonnegative with bounded deterministic row mass before division by
$a$, and its second target is bounded by a fixed multiple of $\zeta$ and
vanishes below the preparatory support threshold.  The buffered row
functions and their witnesses have fixed root Lipschitz constants and
fixed support bounds.  These verify its remaining hypotheses.  The
compression proposition shows that all their errors are a fixed multiple
of $c_*2^{-p}$ when its convergent series of
floors, capacities, and noise weights are chosen.  This includes the
pre-incoming quantities required by the incoming comparison.
Exact endpoint grouping gives common increment arrays for the clamped
and unclamped calculations.  Lemma~\ref{lem:buf-clamping} therefore
preserves these row-function values, including scores and witnesses,
so the same estimates hold after clamping.

Lemma~\ref{lem:atom-assembly} then charges the incoming gates by those
same pre-output norms.  Its separate column-vector calls have a fixed
positive budget slack and decrease task rank.  Equation~\eqref{eq:est-child-precision}
pays their broad-domain factor, and increasing $A_*$ makes their remaining
coefficient as small as required.  A final vector needs no compression;
its clipping is nonexpansive and the bounds at $d_j=1$ apply directly.

Choose the fixed atom range/support bound $C_b$ first.  Choose $c_0$ and
$K_{\rm prec}$ to retain the margin in Equation~\eqref{eq:est-scale-slack}
and fix $B_1$.  The
compression proof then chooses its floor exponent $c_1$, capacity exponent
$c_2$, and required noise decay $K_n$.  Choose task separation and walk
length sufficiently large to supply the needed small $c_*$.  The total
coefficient can thus be made at most $C_b/2$.  Inactive tasks retain their
separate trivial $C_b$ bound.  The remaining half of the active-task
error bound is reserved for the coupled fingerprint error.
\end{proof}

The reserve applies to a local exact-grouping calculation with the stated
averaging property.  The actual algorithm will also use short endpoint
keys.  In Section~\ref{sec:fp-blocks} we define an exact-grouping comparison
within each budget block, verify that it has this averaging property, and
bound its difference from the algorithm.

